\section{Glossaire géométrique-spectral}
\label{app:glossaire}

Le tableau~\ref{tab:glossaire2} récapitule les notations géométriques et spectrales propres à cet article et leur lecture PT.

\begin{table}[htbp]
\centering
\small
\begin{tabular}{@{}p{3.5cm}p{5cm}p{6cm}@{}}
\toprule
Notation & Définition & Lecture PT \\
\midrule
$\Cper$ & Courbe algébrique de persistance dans $\mathcal{KN}$ & Unique solution de ACA$_{\mathcal P}$ avec $|S| \geq 2$ \\
$\genus(\Cper)$ & Genre de $\Cper$ & $\mustar - \lfloor |S|/2 \rfloor = 14$ \\
$\omega_{g,n}^{\Cper}$ & Invariants d'Eynard--Orantin & Cascade des corrélateurs PT \\
$\gPT$ & Métrique de Fisher--Bianchi & Métrique riemannienne canonique sur l'espace des $\mu$ \\
$\Ric(g)$ & Tenseur de Ricci de $g$ & Courbure différentielle de $\gPT$ \\
$\Hess(f)$ & Hessien de $f$ & Hessien du potentiel $f = -\ln \alpha_{\rm sieve}$ \\
$\lambda(\mu)$ & Résidu de Ricci & $-1/\mu^{4} + O(\mu^{-5})$, signature dynamique \\
$H_p(\mu)$ & Taux de Hubble local & Taux de variation de $\thetap{p}$ avec $\mu$ \\
$N(T)$ & Comptage de zéros de Riemann & Loi de Riemann--von~Mangoldt \\
$\Delta_N^{\rm Maslov}$ & Phase de Maslov des coins & Résidu spectral $1/8$ de RvM \\
$\lambdaH$ & Auto-couplage de Higgs & Couplage scalaire à profondeur cascade $d=3$ \\
$\Rres(s)$ & Facteur résiduel zêta régularisé & $\prod_p 1/(1 - \kappa_p(s))$, converge en $\Re(s) > 1$ \\
$\kappa_p(s)$ & Amplitude de canal au premier $p$ & $1 - (1 - p^{-s})\exp(A_p p^{-s})$ \\
$A_p$ & Amplitude PT à $p$ & $\sin^{2}\thetap{p}(\qplus) + \sin^{2}\thetap{p}(\qminus)$ \\
$\HBK$ & Hamiltonien de Berry--Keating régularisé & Opérateur PT $xp$ régularisé \\
\bottomrule
\end{tabular}
\caption{Glossaire géométrique-spectral de l'article 2.}
\label{tab:glossaire2}
\end{table}
